% 4. Formulación de la Función de Evaluación; 5. Análisis Crítico Requerido;
% 7. Descripción de Funciones en el Artículo; 20.1. Criterio 1: Función de Evaluación Othello (20 pts).
\section{Función de Evaluación para Othello}
\textit{Guía: 4. Formulación de la Función de Evaluación; 7. Descripción de Funciones en el Artículo; 20.1. Criterio 1: Función de Evaluación Othello (20 pts).}
\subsection{Formulación y componentes}
Sean $j$ el jugador raíz y $o=3-j$. Para cada color, $N$ cuenta fichas, $C$ esquinas, $M$ movimientos legales y $F$ fichas frontera: ocupadas con al menos un vecino vacío. Definimos $R(a,b)=(a-b)/(a+b)$, con $R(0,0)=0$, y
\begin{equation}
 D=\frac{N_j-N_o}{64},\quad E=\frac{C_j-C_o}{4},\quad
 L=R(M_j,M_o),\quad Q=R(F_o,F_j).
 \label{eq:othello_componentes}
\end{equation}
Los cuatro componentes pertenecen a $[-1,1]$. La frontera aproxima exposición, no estabilidad exacta. Las evaluaciones son
\begin{equation}
 H_{\mathrm{fija}}=10D+40E+35L+15Q,\qquad
 H_{\mathrm{etapas}}=\boldsymbol{w}(t)\cdot(D,E,L,Q),
 \label{eq:othello_dos_funciones}
\end{equation}
con $t=N_j+N_o$. Los pesos aparecen en el Cuadro~\ref{tab:othello_pesos}.
\begin{table}[htbp]
\centering\small
\caption{Pesos de diseño definidos antes del piloto; no ajustados ni demostrados óptimos.}
\label{tab:othello_pesos}
\begin{tabular}{llrrrr}
\toprule
Versión & Condición & $D$ & $E$ & $L$ & $Q$\\
\midrule
Fija & Cualquier $t$ & 10 & 40 & 35 & 15\\
Etapas & $t<20$ & 0 & 35 & 45 & 20\\
Etapas & $20\leq t<45$ & 10 & 40 & 35 & 15\\
Etapas & $t\geq45$ & 55 & 30 & 10 & 5\\
\bottomrule
\end{tabular}
\end{table}
Se priorizan esquinas por su permanencia y movilidad por las opciones disponibles; se penaliza exposición propia y se aumenta el conteo al final. Los umbrales y pesos son hipótesis de ingeniería. Hay al menos tres componentes activos por etapa, pesos en $[0,55]$ y suma 100; por tanto $|H|\leq100$. Esta es la interpretación acordada de ``valores estáticos restringidos'', cuyo intervalo no define la guía. Se cumple $H(s,j)=-H(s,o)$; no se interpreta como probabilidad.

\subsection{Implementación}
Se reutilizan las clases y reglas del docente. Una unidad de profundidad es una ficha colocada; el pase consume cero. MAX/MIN depende del jugador del estado, conservando la perspectiva raíz. La utilidad terminal se comprueba antes del corte:
\begin{equation}
 U(s,j)=\begin{cases}0,&\Delta=0,\\
 \operatorname{sgn}(\Delta)(10000+|\Delta|),&\Delta\neq0,
 \end{cases}\quad\Delta=N_j-N_o.
 \label{eq:othello_terminal}
\end{equation}
Así una victoria domina cualquier hoja heurística. El algoritmo de evaluación cuenta $N,C,M,F$, normaliza según~\eqref{eq:othello_componentes}, selecciona pesos y calcula~\eqref{eq:othello_dos_funciones}. El código íntegro está en \texttt{evaluador.py}. Para un tablero generalizado $n\times n$, conteo y frontera cuestan $O(n^2)$, esquinas $O(1)$ y movilidad $O(n^3)$ por los rayos; ambas evaluaciones cuestan $O(n^3)$ y memoria auxiliar $O(n^2)$. El tablero real es siempre $8\times8$.

\subsection{Análisis crítico}
\textit{Guía: 5. Análisis Crítico Requerido, puntos 1--4.}
\begin{enumerate}
\item Un valor fijo por coordenada ignora líneas y relaciones entre fichas. Maximizar capturas inmediatas puede abrir opciones al rival: la mejora local del conteo no garantiza ventaja final. Nuestras funciones combinan rasgos, pero sus pesos manuales siguen siendo una limitación.
\item La movilidad cuenta alternativas, aunque no distingue su calidad. La guía menciona \emph{sandwich squares} sin definición operativa. Su observación sobre pesos negativos advierte que un patrón aparentemente favorable puede facilitar respuestas rivales; un signo aprendido depende de otros rasgos y del contexto. No se identifica ese término con todas las casillas X/C ni se inventa una penalización. En nuestra función, $Q$ penaliza explícitamente frontera propia; no prueba estabilidad.
\item Los pesos por etapa expresan la transición de flexibilidad a conteo final. Las funciones coinciden en medio juego, por lo que se contrasta el cambio en apertura/final. Los umbrales introducen discontinuidades. Las 13 etapas atribuidas a Logistello en la guía son un antecedente, no una implementación ni resultado nuestro \cite{othello_guia}.
\item La movilidad domina el costo por hoja. Una evaluación costosa puede mejorar decisiones y reducir la profundidad asequible. La cifra de aproximadamente 15\% sobre paridad en Logistello procede de la guía, no de nuestras mediciones. El conteo $D$ no es paridad regional. Como ambas propuestas computan los mismos rasgos, este contraste no aísla el beneficio de añadir cada componente; una ablación futura sería necesaria.
\end{enumerate}
